% Source: https://artofproblemsolving.com/community/c3h2335866_a_problem_about_palindrome_number
% Stable identity: aops-topic-2335866-q01
An $N$-digit palindrome $\overline{d_1d_2\dots d_N}$ has $d_i=d_{N+1-i}$, so it is determined by its first $\left\lceil\dfrac N2\right\rceil$ digits, and these can be chosen freely except that $d_1\ne0$. Hence the number of $N$-digit palindromes is
\[\boxed{9\cdot10^{\left\lceil\dfrac N2\right\rceil-1}},\]
that is, $9\cdot10^{k-1}$ both for $N=2k$ and for $N=2k-1$. For example there are $9$ one-digit, $9$ two-digit, $90$ three-digit and $90$ four-digit palindromes.
